\uuid{Ehkl}
\titre{Résolution d'équation différentielle du premier ordre}
\auteur{bodin ; croizat ; sacre}
\organisation{exo7}
\licence{CC BY-NC-SA 4.0}
\datecreate{2022-03-08}
\niveau{L1}
\module{Analyse}
\chapitre{Equation différentielle}
\sousChapitre{Résolution d'équation différentielle du premier ordre}
\difficulte{2}
\theme{}
\exo7qcmid{hilisit.equations-differentielles.049}

\contenu{
\begin{question}
\qclassification{225.01}
Soit $(E)$ l'équation différentielle $y'+5y = 5x^2 + 2x$. Alors :
\begin{answers}
	\bad{Si $f$ est solution de $(E)$, alors la fonction $x \mapsto f(x)-5x^2-2x$ est solution de l'équation différentielle $(H)$ : $y'+5y = 0$.}
	\good{Si $f$ est solution de $(E)$, alors la fonction $x \mapsto f(x)-x^2$ est solution de l'équation différentielle $(H)$ : $y'+5y = 0$.}
	\bad{Si $f$ est solution de $(E)$, alors la fonction $x \mapsto f(x)-e^{-5x}$ est solution de l'équation différentielle $(H)$ : $y'+5y = 0$.}
	\bad{Si $f$ est solution de $(E)$, alors la fonction $x \mapsto f(x)-2x$ est solution de l'équation différentielle $(H)$ : $y'+5y = 0$.}
\end{answers}
\begin{explanations}
Si $f$ est solution de $(E)$, alors $f'(x) + 5f(x) = 5x^2 + 2x$. On calcule alors que : $(f(x)-x^2)' + 5(f(x)-x^2) = f'(x) - 2x + 5f(x) - 5x^2 = f'(x)+5f(x) - 2x - 5x^2 = 5x^2+2x-2x-5x^2 = 0$. Ainsi la fonction $x \mapsto (f(x)-x^2)$ est bien solution de $(H)$. En revanche, lorsqu'on remplace dans $y'+5y$ avec les fonctions $x \mapsto f(x)-5x^2-2x$, $x \mapsto f(x)-e^{-5x}$ et $x \mapsto f(x)-2x$ (toujours en utilisant le fait que $f'(x)+5f(x)$ peut être remplacé par $5x^2+2x$) on ne trouve pas $0$.
\end{explanations}
\end{question}
}
